\chapter{UAV-based Warehouse Inventory Inspection and QR Code Detection}%
\label{chap:uav-warehouse}

\subimport{.}{intro}

% --- Chapter skeleton (stubs referenced by the roadmap in intro.tex) ---------
% Fill each section below; the \label keys match the \Cref calls in intro.tex.
\section{The Problem of Warehouse Inventory Inspection}%
\label{sec:uav:problem}

This section presents the operational problem that motivates the use of
\glspl{abb:uav}. It first describes the economic and safety stakes of warehouse
inventory inspection, then reviews the main families of methods currently in
use, and finally analyzes their limitations to justify the move toward aerial
platforms.

\subsection{Economic and operational stakes}%
\label{subsec:uav:stakes}

Inventory accuracy is a key driver of warehouse performance. Inaccurate stock
records cause stockouts, overstocking, and order errors, all of which directly
affect customer satisfaction and operating costs~\cite{pore2026drone}. In
traditional warehouses, stock verification relies on periodic stocktaking
cycles in which human operators read rack labels with barcode scanners and
reconcile them against paper logs. This manual process suffers from three
recurring weaknesses:
\begin{itemize}
    \item \textbf{Slow:} a full inventory check in a typical warehouse can take
          more than three days of regular
          operation~\cite{kalinov2020warevision}.
    \item \textbf{Error-prone:} miscounts and misplaced items create
          discrepancies between the \gls{abb:wms} and the actual stock on the
          shelves~\cite{pore2026drone}.
    \item \textbf{Unsafe:} inspecting high racks with forklifts or scissor lifts
          exposes workers to fall hazards~\cite{pore2026drone}.
\end{itemize}

These constraints push warehouses toward more automated solutions. The growth of
e-commerce, the increasing size of storage facilities (both in height and
width), and the demand for real-time stock visibility have made manual
stocktaking economically and operationally difficult to
sustain~\cite{pore2026uavqr}.

\subsection{Traditional and partially automated methods}%
\label{subsec:uav:traditional}

Current inventory inspection methods fall into four families: manual scanning,
\gls{abb:rfid}, ground-based mobile robots, and aerial \glspl{abb:uav}. The
fourth is the focus of this work; the first three---traditional or partially
automated---are summarized below:
\begin{enumerate}
    \item \textbf{Manual scanning:} operators move through the warehouse with
          handheld or fixed barcode scanners to read rack labels and update the
          \gls{abb:wms}~\cite{pore2026drone}. Barcodes, and 1D~barcodes in
          particular, remain the dominant identifier---used in around 87\% of
          warehouses~\cite{kalinov2020warevision}---but the approach is
          low-cost yet labor-intensive, slow, and the main source of human
          errors.
    \item \textbf{\Gls{abb:rfid}:} tags attached to items or pallets allow batch
          reading without a direct line of sight~\cite{pore2026drone}. It speeds
          up operations and removes some manual errors, but performance degrades
          near metal shelving, and the cost of tags and readers limits adoption
          at scale~\cite{pore2026drone}.
    \item \textbf{Ground-based mobile robots (\glspl{abb:agv} and
          \glspl{abb:amr}):} scanners or cameras mounted on a ground platform
          follow predefined or computed routes to inspect
          shelves~\cite{pore2026drone}. They cut labor and improve
          repeatability, but stay confined to the floor plane---limited vertical
          reach and dedicated floor space that competes with picking and
          forklift operations~\cite{pore2026uavqr}.
\end{enumerate}

\subsection{Limits and motivation for UAVs}%
\label{subsec:uav:limits}

Across their differences, these three traditional families share three
structural limitations:
\begin{enumerate}
    \item \textbf{Restricted vertical coverage:} manual scanning depends on
          lifts, and ground robots cannot reach the upper levels of multi-tier
          racks~\cite{pore2026uavqr}.
    \item \textbf{Heavy infrastructure or labor:} each family needs either
          dedicated infrastructure (\gls{abb:rfid} tags and readers,
          \gls{abb:agv}/\gls{abb:amr} fleets, floor markers, charging stations)
          or a significant amount of manual work~\cite{pore2026drone}.
    \item \textbf{Low throughput at scale:} at the scale of a full warehouse
          they remain slow, with stocktaking cycles measured in days rather than
          hours~\cite{kalinov2020warevision}.
\end{enumerate}

\Glspl{abb:uav} are increasingly studied as a complementary technology that
targets these limits directly:
\begin{itemize}
    \item \textbf{Vertical mobility:} they move vertically without lifts, fly
          along narrow aisles, and reach the highest rack levels.
    \item \textbf{Minimal infrastructure:} a charging dock and on-board sensors
          are enough, and missions can run in shifts or at night for
          high-frequency cycle counting without interfering with regular
          operations~\cite{pore2026drone}.
    \item \textbf{Safer operation:} they separate the perception task from the
          human worker, reducing exposure to falls and other
          hazards~\cite{pore2026uavqr}.
    \item \textbf{Lower long-term cost:} initial costs are higher than for
          handheld scanners, but operating costs are lower thanks to reduced
          labor and improved coverage~\cite{pore2026drone}.
\end{itemize}

\subsection{Functional capabilities of a UAV inspection system}%
\label{subsec:uav:capabilities}

Before reviewing the specific challenges of \gls{abb:qr} code detection, swarm
coordination, and simulation covered in the rest of this chapter, it is useful
to clarify what a \gls{abb:uav}-based inspection system actually contains. Such a
system is not a single algorithm but a stack of functional layers that must
operate jointly during a mission:
\begin{itemize}
    \item \textbf{Perception:} relies on cameras and other onboard sensors to
          detect items, identification codes, and obstacles.
    \item \textbf{Localization:} estimates the pose of the \gls{abb:uav} in a
          \gls{abb:gps}-denied environment, from onboard sensing and the
          structural features of the warehouse~\cite{chang2023review}.
    \item \textbf{Planning:} computes the trajectory to follow and, in
          multi-\gls{abb:uav} settings, allocates the inspection workload across
          the fleet.
    \item \textbf{Control:} executes the planned trajectory while keeping the
          \gls{abb:uav} stable in narrow aisles.
    \item \textbf{Communication:} establishes the link between the \gls{abb:uav}
          and a \gls{abb:gcs}, and, when applicable, between cooperating
          \glspl{abb:uav}.
    \item \textbf{Mission management:} coordinates the high-level sequence of
          operations, including takeoff, scanning, docking, and
          recharging~\cite{pore2026drone}.
\end{itemize}

\section{Item Detection and QR Code Recognition from UAVs}%
\label{sec:uav:detection}

Perception is the layer of a \gls{abb:uav}-based inspection system that detects
and decodes the identification codes attached to items and racks. This section
reviews how this layer is implemented in the literature. It first discusses the
choice of \gls{abb:qr} codes among the available identification technologies,
then describes the technical challenges specific to \gls{abb:uav}-based reading,
and finally presents the main families of detection methods.

\subsection{Identification technologies and the choice of QR codes}%
\label{subsec:uav:identification}

Several technologies can be used to identify items in a warehouse. Each presents
different trade-offs in terms of information density, cost, sensing modality, and
robustness:
\begin{itemize}
    \item \textbf{1D~barcodes:} the most widely deployed identifier in industry,
          present in around 87\% of warehouses~\cite{kalinov2020warevision}. They
          are very low cost but encode only a short identifier, are sensitive to
          print quality, and offer no redundancy against damage.
    \item \textbf{\gls{abb:qr} codes:} a 2D matrix code that stores significantly
          more information than a 1D~barcode, includes built-in error correction
          (Reed--Solomon), is cheap to print, and remains readable when partially
          occluded or damaged~\cite{pore2026uavqr, cristiani2020inventory}.
    \item \textbf{\gls{abb:rfid} tags:} allow batch reading without direct line of
          sight, but suffer from signal attenuation near metal racks and require
          dedicated readers and infrastructure~\cite{pore2026drone}.
    \item \textbf{Fiducial markers (ArUco, AprilTag, STag):} provide accurate pose
          estimation when detected by a camera, but encode only an identifier and
          require physical placement of the markers in the
          environment~\cite{do2024experimental}.
\end{itemize}

For \gls{abb:uav}-based warehouse inspection, the \gls{abb:qr} code has emerged as
the most common choice, for three main reasons:
\begin{enumerate}
    \item \textbf{No specialized hardware.} \gls{abb:qr} codes are read with the
          standard onboard cameras already mounted on \glspl{abb:uav} for
          inspection. Some systems reuse the same camera for both navigation and
          inspection, while others add a dedicated camera oriented toward the racks
          for a more favorable viewing angle~\cite{martinez2022warehouse}.
    \item \textbf{Robustness to damage.} The built-in error correction of
          \gls{abb:qr} codes tolerates the partial occlusions and physical damage
          that are typical of warehouse conditions~\cite{pore2026uavqr}.
    \item \textbf{Rich data capacity.} Their large data capacity lets them carry
          not only an identifier but also auxiliary information such as storage
          conditions or product attributes, which simplifies integration with
          \glspl{abb:wms}~\cite{cristiani2020inventory}.
\end{enumerate}
Recent contributions therefore build their entire perception pipeline around
\gls{abb:qr} codes, sometimes complemented by
1D~barcodes~\cite{kalinov2020warevision, alsulami2026deep}.

\subsection{Challenges of QR code detection from a UAV}%
\label{subsec:uav:qr-challenges}

Reading \gls{abb:qr} codes from a \gls{abb:uav} is significantly more difficult
than from a fixed scanner. The literature consistently identifies the following
challenges:
\begin{itemize}
    \item \textbf{Motion blur:} the \gls{abb:uav} is in flight while capturing
          images, which causes blur that degrades the contrast of \gls{abb:qr}
          finder patterns and may prevent
          decoding~\cite{dong2024algorithm, kalinov2020warevision}.
    \item \textbf{Oblique viewing angles:} the \gls{abb:uav} cannot always face
          the \gls{abb:qr} code perpendicularly, especially on multi-tier racks,
          which introduces perspective distortion~\cite{pore2026uavqr}.
    \item \textbf{Variable distance and scale:} as the \gls{abb:uav} moves along
          the aisle or changes altitude, the apparent size of the \gls{abb:qr}
          code changes, requiring scale-robust
          detectors~\cite{kalinov2020warevision, alsulami2026deep}.
    \item \textbf{Uneven warehouse lighting:} artificial lighting, shadows, and
          reflections on metal shelves cause large variations in image contrast
          and can saturate or darken parts of the
          code~\cite{pore2026uavqr, alsulami2026deep}.
    \item \textbf{Limited resolution and small target size:} the \gls{abb:qr} code
          may occupy only a small portion of the image when the \gls{abb:uav}
          scans an entire rack, which makes localization harder and limits the
          available bits for decoding~\cite{kalinov2020warevision}.
    \item \textbf{Occlusion and dense layouts:} packages, neighboring codes, and
          dust on the code can hide part of the pattern, which is particularly
          problematic in densely packed warehouses~\cite{alsulami2026deep}.
\end{itemize}

These challenges have driven the design of dedicated detection methods, which can
be grouped into the four families discussed below.

\subsection{Families of QR code detection methods}%
\label{subsec:uav:qr-families}

Following the taxonomy proposed by Alsulami and
Jhanjhi~\cite{alsulami2026deep}, \gls{abb:qr} code detection methods used in
\gls{abb:uav}-based inventory can be organized into four families.

\begin{enumerate}[itemsep=0.7em]
    \item \textbf{Classical image-processing approaches.} These methods rely on
          hand-crafted features such as edge detection, gradient analysis, and
          pattern matching to locate the \gls{abb:qr} finder patterns, followed by
          standard decoding libraries (ZBar, Pyzbar, OpenCV's \gls{abb:qr}
          decoder). They are fast, lightweight, and easy to deploy onboard a
          \gls{abb:uav}, with detection rates around 90\% under controlled flight
          conditions~\cite{martinez2022warehouse, cristiani2020inventory}.
          However, their performance degrades sharply under motion blur, low
          light, or oblique angles, which restricts them to slow, well-lit, and
          stable flights~\cite{alsulami2026deep}.

    \item \textbf{Deep learning--based detection.} \Glspl{abb:cnn} and one-stage
          object detectors such as YOLO have become the dominant approach for code
          localization in \gls{abb:uav} imagery, as they learn directly from data
          and are more robust to perspective distortion, lighting variations, and
          small target sizes. A recent benchmark on a simulated \gls{abb:uav}
          barcode dataset reports an mAP@0.5 of 92.4\% with YOLOv8, at inference
          times below 4~ms per image---compatible with real-time onboard
          operation~\cite{alsulami2026deep}. \gls{abb:cnn}-based detectors have
          also been embedded in production stocktaking systems, where they correct
          the flight trajectory from detected codes to improve the recognition
          rate over fixed scanning patterns~\cite{kalinov2020warevision}.

    \item \textbf{Hybrid approaches.} These combine deep-learning localization with
          classical decoding: the network only locates and crops the \gls{abb:qr}
          region, while a standard decoder (e.g.\ OpenCV's \gls{abb:qr} module)
          extracts the data. This keeps localization robust while decoding stays
          simple and standards-compliant. Reported accuracies are high in good
          conditions (around 98\% when well-lit) but drop under occlusion (around
          78\%) or low light (around 85\%)~\cite{alsulami2026deep}.

    \item \textbf{\gls{abb:uav}-specific pipelines.} This last family targets
          degradations that are specific to \gls{abb:uav} imagery. One approach
          uses a \gls{abb:gan} with attention mechanisms to deblur \gls{abb:qr}
          images before decoding, recovering codes that classical methods fail to
          read~\cite{dong2024algorithm}. Another is \emph{active perception}, where
          the \gls{abb:uav} adapts its trajectory from the codes already detected
          to improve subsequent observations; this raises the recognition rate from
          around 73--96\% with fixed overlapping-grid trajectories to over 99\% in
          a controlled lab warehouse setup~\cite{kalinov2020warevision}.
\end{enumerate}

A synthesis of representative \gls{abb:qr} detection works in the \gls{abb:uav}
inventory context is given in \Cref{tab:uav:qr-detection}.

\begin{table}[htbp]
    \centering
    \caption{Representative QR / barcode detection methods for UAV-based
        inventory.}
    \label{tab:uav:qr-detection}
    \small
    \begin{tabular}{@{}l
            >{\raggedright\arraybackslash}p{2.5cm}
            >{\raggedright\arraybackslash}p{3.6cm}
            >{\raggedright\arraybackslash}p{2.7cm}
            >{\raggedright\arraybackslash}p{3.3cm}@{}}
        \toprule
        \textbf{Ref} & \textbf{Family} & \textbf{Technique} &
        \textbf{Validation} & \textbf{Result} \\
        \midrule
        \cite{martinez2022warehouse} & Classical & OpenCV QR decoder & Mock warehouse & $\sim$90\% detection \\
        \cite{cristiani2020inventory} & Classical & Standard QR library & Lab testbed & Speed/accuracy trade-off \\
        \cite{kalinov2020warevision} & DL + active perception & CNN segmentation + trajectory correction & Lab testbed & 99.3\% vs.\ 73--96\% (fixed) \\
        \cite{alsulami2026deep} & DL (hybrid) & YOLOv8 + OpenCV decoder & Images only & mAP@0.5 = 92.4\% \\
        \cite{dong2024algorithm} & UAV-specific & GAN + attention (deblurring) & Images only & High under motion blur \\
        \cite{pore2026uavqr} & DL + safety & Multi-UAV QR scanning & Lab testbed + 3D simulation & 95--96\% sim / 86--90\% real \\
        \bottomrule
    \end{tabular}

    \vspace{0.8em}
    \begin{minipage}{0.95\linewidth}
        \footnotesize
        \textbf{Validation levels:}
        \par\vspace{2pt}
        \renewcommand{\arraystretch}{1.15}
        \begin{tabular}{@{}>{\itshape}l @{\hspace{0.8em}} p{0.66\linewidth}@{}}
            Images only    & static image dataset; no physical drone or simulator.        \\
            3D simulation  & virtual drone flying in a simulated 3D environment.           \\
            Lab testbed    & real drone in a small laboratory setup.                       \\
            Mock warehouse & real drone in a setup that physically reproduces a warehouse. \\
        \end{tabular}
    \end{minipage}
\end{table}

% Garde le tableau ci-dessus dans la section 1.3 : il doit être placé avant
% que la section 1.4 ne commence (sinon il « coupe » la nouvelle section).
\FloatBarrier
\section{From Single-UAV to Multi-UAV and Swarm Coordination}%
\label{sec:uav:swarm}

Most of the \gls{abb:uav}-based inventory systems presented so far rely on a
single \gls{abb:uav}. This section first discusses the limits of
single-\gls{abb:uav} approaches and the motivation for moving toward
multi-\gls{abb:uav} systems, then reviews the main families of coordination
architectures reported in the literature, together with their typical
communication and fault-tolerance properties.

\subsection{Limits of single-UAV approaches and motivations for swarms}%
\label{subsec:uav:single-limits}

Using a single \gls{abb:uav} for warehouse inventory inspection faces several
intrinsic limitations:
\begin{itemize}[itemsep=0.4em]
    \item \textbf{Limited coverage.} A single \gls{abb:uav} must traverse the whole
          warehouse one aisle at a time, so the mission duration grows linearly
          with the size of the facility and the number of
          racks~\cite{pore2026drone}.
    \item \textbf{Limited battery autonomy.} Indoor \glspl{abb:uav} typically fly
          for only 10 to 30 minutes per charge, so a single platform may need
          several recharging cycles to inspect a large facility, which fragments
          the mission and limits real-time stock
          visibility~\cite{pore2026drone, kumar2025comprehensive}.
    \item \textbf{Single point of failure.} If the only \gls{abb:uav} fails
          (battery depletion, sensor fault, collision), the mission is interrupted
          and must be restarted---a critical issue in time-constrained
          operations~\cite{chandran2024network}.
    \item \textbf{Lack of redundancy.} With a single \gls{abb:uav}, a missed
          \gls{abb:qr} code or an undetected obstacle cannot be cross-checked from
          another viewpoint.
\end{itemize}

These limits motivate multi-\gls{abb:uav} systems, in which several drones
cooperate during the same mission. The benefits commonly cited in the literature
are:
\begin{itemize}[itemsep=0.4em]
    \item \textbf{Parallelization} of the workload, which reduces the total mission
          time by distributing aisles or racks across
          \glspl{abb:uav}~\cite{pore2026uavqr, chandran2024network}.
    \item \textbf{Redundancy}, since several \glspl{abb:uav} can observe the same
          item from different viewpoints, improving detection robustness.
    \item \textbf{Scalability}, as inspection capacity grows by adding
          \glspl{abb:uav} rather than upgrading a single
          platform~\cite{kumar2025comprehensive}.
    \item \textbf{Fault tolerance}, since the failure of one \gls{abb:uav} can be
          absorbed by reassigning its tasks to others without interrupting the
          mission~\cite{chandran2024network}.
\end{itemize}

\subsection{Families of coordination architectures}%
\label{subsec:uav:architectures}

Coordinating several \glspl{abb:uav} raises the question of how they share
information and make decisions. The literature distinguishes three main families
of architectures:
\begin{enumerate}[itemsep=0.7em]
    \item \textbf{Centralized architectures.} All \glspl{abb:uav} communicate with
          a central \gls{abb:gcs} that performs global path planning, task
          allocation, and conflict resolution. This is the most common
          architecture in practice, as it is simple to implement and minimizes
          onboard computation; its star-shaped topology, however, adds
          communication overhead and creates a single point of failure at the
          \gls{abb:gcs}~\cite{kumar2025comprehensive}. A representative example is
          the multi-\gls{abb:uav} system of Pore et~al.~\cite{pore2026uavqr}, where
          two \glspl{abb:uav} receive sector assignments from a \gls{abb:gcs} that
          runs the safety-aware trajectory planner and aggregates the \gls{abb:qr}
          scanning results.
    \item \textbf{Distributed architectures.} Each \gls{abb:uav} performs part of
          the decision-making locally, from information exchanged with its
          neighbors, while the \gls{abb:gcs}---when present---only handles
          initialization or high-level supervision. This relieves the central node
          and improves resilience to communication failures, at the cost of more
          complex onboard algorithms and increased peer-to-peer
          communication~\cite{kumar2025comprehensive}. A typical scheme is
          \emph{leader--follower}, where one \gls{abb:uav} coordinates while the
          others follow~\cite{chandran2024network}.
    \item \textbf{Decentralized architectures.} This family removes central control
          entirely: each \gls{abb:uav} independently gathers and shares information
          about its location and environment, and collaborative decisions emerge
          from local interactions~\cite{kumar2025comprehensive}. It is especially
          advantageous when central control is impractical or when high autonomy is
          required. Chandran and Vipin~\cite{chandran2024network} propose such a
          strategy, in which \glspl{abb:uav} self-organize into a linear formation,
          optimize their coverage paths, and adapt to agent failures, letting a
          swarm of up to ten \glspl{abb:uav} keep its formation despite node
          failures.
\end{enumerate}

Across these three families, two cross-cutting concerns strongly affect
performance. The first is the \emph{communication model}: in real warehouses the
wireless channel is noisy, bandwidth is limited, and packet losses are frequent.
Recent work evaluates coordination strategies under realistic network conditions
with the NS-3 network simulator, measuring metrics such as packet-delivery ratio,
throughput, and delay~\cite{chandran2024network}. The second is \emph{fault
tolerance}: the ability of the swarm to keep going when one or several
\glspl{abb:uav} fail. Chandran and Vipin~\cite{chandran2024network}, for instance,
re-elect a leader \gls{abb:uav} when the current one becomes unavailable, so the
swarm completes its coverage path despite node failures.


\section{Simulators and Co-Simulation Frameworks}%
\label{sec:uav:evaluation}

Validating a \gls{abb:uav}-based inspection system before any field deployment
relies on simulation environments that reproduce both the physical behavior of
the drone and the communication network between drones. This section reviews the
main \gls{abb:uav} simulators reported in the literature, then the co-simulation
frameworks that couple them with network simulators to make multi-\gls{abb:uav}
validation realistic.

\subsection{Simulators for UAV systems}%
\label{subsec:uav:simulators}

A \gls{abb:uav} simulator is a software environment that reproduces the behavior
of a drone and the data produced by its onboard sensors, so that the same
control, perception, and navigation algorithms that would run on a real
\gls{abb:uav} can be developed and tested entirely in a computer. Around this
central role, a few layers have become standard across the research community:
\begin{itemize}[itemsep=0.4em]
    \item Most simulators expose interfaces compatible with the \gls{abb:ros}, a
          middleware that connects the software modules of a robotic system.
    \item They run the autopilot---typically the open-source PX4 or
          ArduPilot---in \gls{abb:sitl} mode, where the real autopilot code is
          executed against the simulated drone rather than on physical flight
          hardware, the two communicating through the MAVLink protocol.
    \item The \emph{photo-realism} of the rendered scene (the visual fidelity of
          camera images, lighting, and materials) has become a key differentiator
          between modern frameworks, since any perception algorithm that processes
          visual data depends on the realism of these images.
\end{itemize}

Within this landscape, the Pegasus Simulator paper provides a recent and
well-known comparison of the main frameworks used in \gls{abb:uav}
research~\cite{jacinto2024pegasus}. The most widely used ones are:
\begin{itemize}[itemsep=0.4em]
    \item \textbf{Gazebo} is the most established simulator in the robotics
          community. Tightly integrated with \gls{abb:ros}, it offers a rigid-body
          physics engine and a PX4-\gls{abb:sitl} extension for autopilot
          integration; its main limitation is its OpenGL-based rendering, which
          lacks the visual fidelity of modern game
          engines~\cite{jacinto2024pegasus, acharya2020cornet}.
    \item \textbf{AirSim}, developed by Microsoft, is built on Unreal Engine (or
          Unity), producing photo-realistic graphics. It provides rich sensor
          models, PX4 integration, and a wrapper for OpenAI~Gym (a standard
          interface for training reinforcement-learning agents), and underpins
          several recent warehouse-oriented
          pipelines~\cite{jacinto2024pegasus, tejero2024realistic}.
    \item \textbf{Flightmare}, developed at UZH, uses a custom physics engine and
          Unity for rendering. It offers fast simulation and OpenAI~Gym
          integration, but no out-of-the-box PX4
          support~\cite{jacinto2024pegasus}.
    \item \textbf{jMAVSim and MuJoCo} are lighter alternatives: jMAVSim is a
          Java-based PX4 community simulator without photo-realistic rendering,
          used mainly to test basic autopilot functions, while MuJoCo is a
          general-purpose physics engine widely used in control and
          reinforcement-learning research, with limited integration with flight
          stacks~\cite{jacinto2024pegasus}.
    \item \textbf{Pegasus Simulator} is a recent extension built on NVIDIA Isaac
          Sim, combining a high-quality photo-realistic renderer (RTX) with a
          modern physics engine (PhysX). Pegasus adds the PX4 and ROS\,2
          integration that was previously missing, and exposes a Python API for
          fast prototyping of guidance, navigation, and control algorithms on
          multiple aerial vehicles~\cite{jacinto2024pegasus}.
\end{itemize}

Beyond these generic frameworks, several recent works build dedicated simulation
pipelines for indoor warehouse inspection. Tejero-Ruiz et~al.\ propose a
comprehensive architecture combining AirSim, Unreal Engine, and \gls{abb:ros},
with Docker-based deployment (which packages software and its dependencies into
isolated containers so that the same simulation is reproducible on different
machines), ArduPilot/PX4 \gls{abb:sitl}, realistic \gls{abb:lidar} noise models,
and configurable static and dynamic objects~\cite{tejero2024realistic}. The
\gls{abb:qr} scanning system of Pore et~al.\ instead uses a custom co-simulation
between MATLAB--Simulink (an engineering tool for modeling dynamical systems) and
Unreal Engine, with photo-realistic warehouse scenes populated with
\gls{abb:qr}-labeled boxes and multi-tier racks~\cite{pore2026uavqr}. These
tailored pipelines indicate that, despite the maturity of generic simulators, the
warehouse scenario is not natively covered by mainstream frameworks and still
requires domain-specific extensions.

A comparison of the main \gls{abb:uav} simulators along their core features is
given in \Cref{tab:uav:simulators}.

\begin{table}[htbp]
    \centering
    \caption{Comparison of the main UAV simulators, adapted from Table~I
        in~\cite{jacinto2024pegasus}.}
    \label{tab:uav:simulators}
    \small
    \begin{tabular}{@{}l >{\raggedright\arraybackslash}p{2.6cm} c c c c c@{}}
        \toprule
        \textbf{Platform} & \textbf{Base engine} & \textbf{Photo} & \textbf{Sensors} & \textbf{Vision} & \textbf{MAVLink} & \textbf{API} \\
        \midrule
        jMAVSim    & Java                     &            & \checkmark &            & \checkmark & $\bigstar$ \\
        RotorS     & Gazebo                   &            & \checkmark & \checkmark & \checkmark & $\bigstar\bigstar$ \\
        PX4-SITL   & Gazebo                   &            & \checkmark & \checkmark & \checkmark & $\bigstar\bigstar$ \\
        AirSim     & Unreal Engine (or Unity) & \checkmark & \checkmark & \checkmark & \checkmark & $\bigstar\bigstar\bigstar$ \\
        Flightmare & Unity                    & \checkmark & \checkmark & \checkmark &            & $\bigstar\bigstar\bigstar$ \\
        MuJoCo     & OpenGL (or Unity)        &            &            &            &            & $\bigstar\bigstar$ \\
        Pegasus    & NVIDIA Isaac Sim         & \checkmark & \checkmark & \checkmark & \checkmark & $\bigstar\bigstar$ \\
        \bottomrule
    \end{tabular}

    \vspace{1em}
    \begin{minipage}{0.95\linewidth}
        \footnotesize
        \textbf{Legend.}\quad \checkmark{} = feature present;\quad empty cell = feature absent.
        \par\vspace{5pt}
        \renewcommand{\arraystretch}{1.3}
        \begin{tabular}{@{}>{\bfseries}l @{\hspace{1.2em}} >{\raggedright\arraybackslash}p{0.7\linewidth}@{}}
            Photo   & photo-realistic rendering                                            \\
            Sensors & simulated onboard IMU\,/\,GPS                                         \\
            Vision  & camera (vision) sensors                                              \\
            MAVLink & MAVLink\,/\,PX4 interface                                             \\
            API     & API complexity, from $\bigstar$ (simplest) to $\bigstar\bigstar\bigstar$ (most complex) \\
        \end{tabular}
    \end{minipage}
\end{table}

\FloatBarrier
\subsection{Network simulation and co-simulation for swarm validation}%
\label{subsec:uav:cosim}

A network simulator is a software environment that reproduces the behavior of a
communication network---packet transmissions, delays, losses, and routing
protocols---for technologies such as Wi-Fi, LTE, or 5G. When an inspection task
involves several drones, the wireless links between them, and between them and
the \gls{abb:gcs}, become a critical part of the system, since communication
delays and packet losses directly affect coordination and safety. Realistic
validation of multi-\gls{abb:uav} systems therefore requires not only a physics
simulator but also a network simulator. Two open-source network simulators
dominate the literature: NS-3 and OMNeT++~\cite{acharya2020cornet}. NS-3 in
particular is widely used in \gls{abb:uav} swarm research, thanks to its support
for both wired and wireless networks and its modular library, which eases
integration with external tools such as physics
simulators~\cite{chandran2024network}.

The combination of a \gls{abb:uav} physics simulator with a network simulator is
called a \emph{co-simulation framework}. The two simulators run in parallel and
exchange information at each time step: the physics simulator provides the
position and motion of each \gls{abb:uav}, while the network simulator computes
the resulting communication delays and packet losses, which then feed back into
the control loop. Such setups are essential for testing swarm coordination under
realistic communication conditions. Several have been proposed:
\begin{itemize}[itemsep=0.4em]
    \item \textbf{CUSCUS} integrates the FL-AIR flight simulator with NS-3 through
          tap-bridge containers; its main limitation is the difficulty of
          extending it to LTE or 5G~\cite{acharya2020cornet}.
    \item \textbf{AVENS} combines OMNeT++ with the X-Plane flight simulator, but
          the two only exchange \gls{abb:uav} positions through an XML file, with
          no path for control or telemetry~\cite{acharya2020cornet}.
    \item \textbf{FlyNetSim} is built on NS-3 and ArduPilot \gls{abb:sitl}, with
          timestamped packets for synchronization~\cite{acharya2020cornet}.
    \item \textbf{CORNET} connects Gazebo (\gls{abb:uav} physics) with NS-3
          (network) through a \gls{abb:ros}-based middleware. It provides time
          synchronization between the two simulators, treats late packets as
          losses, and is open source~\cite{acharya2020cornet}.
\end{itemize}

The decentralized swarm coordination work of Chandran and
Vipin~\cite{chandran2024network} is a representative use case: their strategy is
validated in two phases---mission efficiency and fault tolerance on \gls{abb:ros}
and Gazebo, and inter-\gls{abb:uav} communication on NS-3. This kind of two-phase,
co-simulated validation is increasingly seen as a prerequisite for credible swarm
research.

\section{Chapter Conclusion}%
\label{sec:uav:conclusion}

The state of the art reviewed in this chapter reveals a clear contrast between
layers. The \emph{perception} and \emph{platform} layers of \gls{abb:uav}-based
warehouse inspection are mature: \gls{abb:qr} code detectors reach high
accuracy~\cite{dong2024algorithm, alsulami2026deep}, onboard sensing enables
flight in \gls{abb:gps}-denied environments~\cite{chang2023review}, and several
\gls{abb:uav}-based inspection systems have been demonstrated in laboratory and
mock-warehouse setups~\cite{kalinov2020warevision, pore2026uavqr, cristiani2020inventory, martinez2022warehouse}.
The \emph{coordination}, \emph{decision-making}, and \emph{validation} layers, by
contrast, are still developing, and four main gaps stand out across the
literature:
\begin{enumerate}[itemsep=0.7em]
    \item \textbf{Scale of inspection.} Around 83\% of warehouse-drone studies
          still rely on a single \gls{abb:uav}~\cite{pore2026drone}, and
          multi-\gls{abb:uav} coordination in narrow indoor aisles is reported in
          only a small fraction of works.
    \item \textbf{Simulation environments.} Despite the maturity of
          general-purpose \gls{abb:uav} simulators, the indoor warehouse scenario
          is not natively supported by mainstream tools and typically requires
          custom extensions on top of frameworks such as AirSim or
          Simulink--Unreal~\cite{pore2026uavqr, tejero2024realistic}, which makes
          results difficult to compare across studies.
    \item \textbf{Communication realism.} Although the wireless links between
          \glspl{abb:uav} strongly affect coordination, only a minority of works
          integrate a network simulator into their validation
          pipeline~\cite{chandran2024network, acharya2020cornet}, so many reported
          multi-\gls{abb:uav} results implicitly assume idealized communication.
    \item \textbf{Dynamic conditions.} Most studies still operate in static
          environments, while dynamic warehouse conditions---human operators,
          forklifts, evolving stock---remain comparatively
          under-explored~\cite{pore2026drone}, which limits the transferability of
          reported results to industrial deployments.
\end{enumerate}

The contribution of this thesis is positioned at the intersection of these four
gaps. It proposes a multi-\gls{abb:uav} decision-making approach (gap~1),
validated in a dynamic warehouse environment built on a modern simulator
(gaps~2 and~4), with explicit modeling of the wireless communication through
robot--network co-simulation (gap~3). The next chapter introduces the
methodological foundations of the proposed approach and details how each of these
gaps is addressed.
